\documentclass[11pt]{article}
\usepackage{amsmath,amsthm,mathtools,tasks,tcolorbox,multicol,siunitx,fancyhdr,lastpage}
\pagestyle{fancy}
\fancyhf{}
\lhead{Trường THPT Lê Lợi}
\rhead{Đề số 01}
\cfoot{Trang \thepage/\pageref{LastPage}}
\definecolor{xanh}{RGB}{0,102,204}
\newcommand*{\mycmd}[2][Mặc định]{\textbf{#1}: #2}
\NewDocumentCommand{\diemso}{s m O{điểm}}{\IfBooleanTF{#1}{(#2 #3)}{[#2 #3]}}
\DeclarePairedDelimiter{\ceilx}{\lceil}{\rceil}
\def\vect#1{\overrightarrow{#1}}
\def\khoang#1;#2;{\left[#1;#2\right]}
\let\epsi\varepsilon
\theoremstyle{definition}
\newtheorem{dn}{Định nghĩa}[section]
\newtheorem*{nx}{Nhận xét}
\begin{document}
 
\section{Các macro lạ}
\textbf{Mặc định}: giá trị và \textbf{Tùy chọn}: khác. Điểm: [2 điểm] hoặc (3 đ).
Làm tròn $\left\lceil x \right\rceil$, vectơ $\overrightarrow{MN}$, khoảng $\left[a;b\right]$, $\varepsilon>0$.
 Văn bản trong lệnh lạ  rồi  \ tiếp tục  2cm nội dung giữ lại .



\textbf{Định nghĩa 1.} 
Hàm số $f$ liên tục tại $x_0$ nếu $\lim\limits_{x\to x_0} f(x)=f(x_0)$.





\textbf{Nhận xét 1.} 
Theo Định nghĩa~[dn:1], \textit{mọi đa thức} đều liên tục.





\textbf{Chứng minh.} 
Hiển nhiên.




\subsection{Toán học nâng cao}
\begin{equation}\label{eq:main}
\sum_{k=1}^{n} \binom{n}{k} = 2^n - 1, \quad \operatorname*{argmax}_{x\in X} g(x)
\end{equation}
\begin{gather*}
A = \begin{pmatrix} 1 & 2 \\ 3 & 4 \end{pmatrix}, \quad \det A = \begin{vmatrix} 1 & 2\\ 3 & 4\end{vmatrix} = -2
\end{gather*}
\[
f(x)=\begin{cases} x^2 & \text{nếu } x\ge 0\\ -x & \text{nếu } x<0 \end{cases}, \qquad
\underbrace{1+1+\cdots+1}_{n \text{ lần}} = n, \qquad \mathcal{F}, \mathfrak{g}, \mathbb{Z}_p
\]
Phương trình ([eq:main]) và $x \equiv 1 \pmod 3$, $a \nmid b$, $\vec{u}\perp\vec{v}$, $\widehat{ABC}=60^\circ$.

\subsection{Danh sách và bố cục}


a) $x=1$

b) $x=2$

c) $x=3$

d) $x=4$




\textbf{Định lí Vi-ét} Tổng hai nghiệm $x_1+x_2=-\frac{b}{a}$.

\textbf{Hệ quả} Tích hai nghiệm.




1. Mức 1
  

(a) Mức 2
    

i. Mức 3

2. Mức 1 tiếp





\textbf{Ghi nhớ}


Công thức nghiệm: $x=\dfrac{-b\pm\sqrt{\Delta}}{2a}$.







Cột trái với nhiều chữ để thấy hai cột được chia ra trong Word, có đủ nội dung để kiểm tra.
 
Cột phải cũng có nội dung.







\textit{Người ra đề}\\
\textbf{Nguyễn Văn A}




\begin{figure}[ht]
\centering


ZLZIMGZ00001ZEND


\caption{Đồ thị hàm số}
\end{figure}
Hình~[fig:1] cho thấy\ldots{} Trích dẫn  sgk .

\begin{tabularx}{\textwidth}{|X|c|}
\hline
  Cột rộng & Hẹp \\ \hline
Nội dung dài dài dài &   1 \\ \hline
\end{tabularx}




 sgk  Bộ GD\&ĐT, \emph{Sách giáo khoa Toán 12}, NXB Giáo dục, 2024.




Ký tự đặc biệt: \% \$ \& \# \_ \{ \} \textbackslash{}  \    50\,\% --- dấu ``nháy'' và ~cách. Accent: \'e \`a \^o \"u \~n  c   a   a .
Mã ZLZVERBZ00001ZEND giữ nguyên.
\end{document}
